\subsection*{Procesamiento}
% Dar formato a toda la información sin filtrar obtenida en la fase anterior. En caso de querer obtener más información explicar el posible uso de la misma, así como su relevancia.

Filtramos la información cruda de las herramientas limpiándola del ruido y comentarios siguiendo las siguientes estrategias para facilitar la lectura e interpretación del resultado:

\begin{itemizeColor}
    \item \textbf{Limpieza de texto:} En las consultas de WHOIS filtramos mediante expresiones regulares los avisos de privacidad y licencias para quedarnos solo con los datos de contacto, fechas y servidores de nombres. De igual forma, en las IPs extraemos exclusivamente los bloques CIDR y rangos de red asignados. Además, en las pruebas de conectividad con \texttt{ping} filtramos el texto completo para aislar directamente el porcentaje de paquetes perdidos y las estadísticas de latencia (mínima, media y máxima).
    \item \textbf{Normalización y verificación de subdominios:} Al combinar herramientas se generan algunos resultados repetidos. El script los unifica, elimina duplicados con \texttt{sort -u}, descarta aquellos que ya no resuelven a una IP activa mediante \texttt{dig} y les da formato de tabla. Además, con la bandera \texttt{-n} limitamos la cantidad de filas mostradas para que el reporte no se vuelva excesivamente largo, o podemos prescindir completamente de este paso mediante \texttt{-s} en dominios muy extensos.
    \item \textbf{Aislamiento de puertos útiles:} En Nmap descartamos los puertos cerrados o filtrados (\texttt{--open}), agrupando los resultados por IP con sus versiones de servicio identificadas.
\end{itemizeColor}
